\section{Vocabulario y conceptos clave}
Una parte esencial del enfoque CPRE es trabajar con terminolog\'ia precisa y compartida. El handbook define conceptos base que deben utilizarse de forma consistente en toda especificaci\'on \cite[pp.~10--11]{cpre}.

\begin{itemize}[leftmargin=1.2em]
  \item \textbf{Requirement:} necesidad percibida, capacidad del sistema o representaci\'on documentada de ambas.
  \item \textbf{Requirements specification:} colecci\'on sistem\'atica de requisitos que cumple criterios de calidad.
  \item \textbf{Requirements Engineering:} enfoque sistem\'atico y disciplinado para especificar y gestionar requisitos.
  \item \textbf{Stakeholder:} persona u organizaci\'on que influye en requisitos o es impactada por el sistema.
  \item \textbf{System / context / scope:} delimitaciones clave para entender qu\'e se dise\~na, qu\'e se asume y qu\'e queda fuera.
\end{itemize}

El cap\'itulo de glosarios del documento refuerza que la terminolog\'ia acordada debe mantenerse de forma centralizada, con uso obligatorio en los artefactos del proyecto para evitar ambig\"uedades y malentendidos \cite[pp.~72--73]{cpre}.
